\input{welibc}

% Babel Publishing journal skin for this paper.
\definecolor{BabelInk}{HTML}{1C1A17}
\definecolor{BabelRule}{HTML}{5F3B22}
\definecolor{BabelSoft}{HTML}{F2EFE8}
\definecolor{BabelAccent}{HTML}{8A1F2D}
\definecolor{BabelTableHead}{HTML}{E7DFD2}

% This paper uses \hebrew{...} for every Hebrew word.
% Default: regular square Hebrew.
\newcommand{\hebrew}[1]{\heb{#1}}
% To print every Hebrew word in Paleo-Hebrew instead, uncomment this line:
% \renewcommand{\hebrew}[1]{\paleo{#1}}

\newcommand{\JournalName}{Babel Publishing}
\newcommand{\JournalIssue}{Journal of Scriptural Computation 1.1}
\newcommand{\JournalDOI}{doi:10.0000/babel.four-streams}

\newcommand{\BabelLogo}{%
  \begingroup
  \setlength{\fboxsep}{6pt}%
  \setlength{\fboxrule}{0.55pt}%
  \fcolorbox{BabelRule}{BabelSoft}{%
    \begin{minipage}{0.31\linewidth}
    \centering\small\linespread{0.92}\selectfont
    \chineseA{肏}IV\textbar{}\textbar{}\chineseA{一}\\
    Fʌ\korean{ㅋ}\ara{𐡅}v∨\paleo{י}\\
    φə\paleo{ק}\chineseA{愛}\heb{ו}\textbar{} \heb{י}
    \end{minipage}%
  }%
  \endgroup
}

\newcommand{\JournalTable}[2]{%
  \par\smallskip
  \begin{center}
  \small
  \renewcommand{\arraystretch}{1.18}%
  \setlength{\tabcolsep}{4.8pt}%
  \rowcolors{2}{BabelSoft}{white}%
  \arrayrulecolor{BabelRule}%
  \Table{#1}{\hline\rowcolor{BabelTableHead}#2}%
  \arrayrulecolor{black}%
  \rowcolors{2}{}{}%
  \end{center}
  \smallskip
}

\newenvironment{JournalAbstract}{%
  \section*{Abstract}%
  \begingroup
  \small
  \leftskip=1.8em
  \rightskip=1.8em
  \noindent\ignorespaces
}{\par\endgroup\medskip}

\newcommand{\JournalQuote}[1]{%
  \begin{quote}\small\itshape #1\end{quote}%
}

\newcommand{\ArticleTitleBlock}{%
  \thispagestyle{empty}%
  \noindent\begin{minipage}[t]{0.61\linewidth}
    {\footnotesize\scshape\color{BabelRule}\JournalName\par}
    {\scriptsize\color{BabelInk}\JournalIssue\quad\JournalDOI\par}
  \end{minipage}\hfill\BabelLogo
  \vspace{2.2em}

  \begin{center}
    {\Huge\bfseries\color{BabelInk} The Four Streams\par}
    \vspace{0.5em}
    {\large\itshape The Nameless: Genesis 2:10--14, word boundaries, and a testable second reading\par}
    \vspace{1.1em}
    {\normalsize Anonymous Research Note\par}
    \vspace{0.4em}
    {\small Received August 2026; published by Babel Publishing\par}
  \end{center}
  \vspace{1.2em}
  \noindent\color{BabelRule}\rule{\linewidth}{0.7pt}\color{BabelInk}
  \vspace{0.8em}
}

\begin{document}
\pagenumbering{arabic}
\pagestyle{fancy}
\fancyhf{}
\fancyhead[LE,RO]{\small\thepage}
\fancyhead[LO]{\small\itshape The Four Streams}
\fancyhead[RE]{\small\scshape Babel Publishing}
\renewcommand{\headrulewidth}{0.3pt}
\renewcommand{\footrulewidth}{0pt}
\setlength{\parindent}{1.1em}
\setlength{\parskip}{0.25em}
\color{BabelInk}

\def\CurrentBook{Babel Publishing}%
\def\CurrentTitleBook{Babel Publishing}%
\def\CurrentApocryphaHeader{Babel Publishing}%
\def\CurrentChapter{The Four Streams}%
\SetCurrentBookImageDir{Babel Publishing}%

\ArticleTitleBlock
\begin{JournalAbstract}

Genesis 2:10 says something almost too convenient for this paper: one river leaves Eden, then divides and becomes four ``heads.'' The ordinary meaning is geographical. This paper asks whether the writing may also perform, at the level of consonants, something like the thing it describes.

The proposal is narrow. In the Masoretic consonantal tradition, a run beginning with the final two words of Genesis 2:11 and continuing through all of 2:12 can be divided into words a second way without inserting, deleting, replacing, or rearranging a consonant. One consonant left hanging at the end of 2:12 resolves immediately when the first word of 2:13 is included. I call this a \textbf{zero-edit resegmentation}: same letters, same order, different word boundaries. The ordinary reading remains ordinary; the question is whether a second lexical structure was also made available.

The result is more interesting---and less sensational---than a ``hidden message.'' Most of the alternate pieces are straightforward Hebrew. Particularly strong are a phrase meaning roughly ``with/in her I will be pleased,'' an ordinary ``he is good,'' the rare but exact Biblical Hebrew noun ``a severed piece,'' the ordinary noun ``son/offspring,'' and a word for an individual sheep-or-goat. A recurring three-consonant packet can independently mean ``there,'' ``name,'' or ``he placed.'' The weakest piece is a three-letter packet that does \textbf{not} have a secure ordinary Biblical Hebrew interpretation adequate to the proposed sentence. That weakness matters. It prevents us from pretending that the experiment has produced a polished secret paragraph.

The strongest conclusion is therefore structural, not theological: the Masoretic consonants admit an unusually suggestive second segmentation whose lexical density and local semantic echoes deserve measurement against controls. The strongest objection is textual: the Samaritan Pentateuch has a materially different Genesis 2:12, including an extra ``very'' and a different consonantal spelling of the feminine pronoun, so the pattern is a property of the Masoretic textual branch as we possess it, not something we may casually project back to an inaccessible autograph.\cite{MasoreticGenesis,SamaritanGenesis,Naaijer2024,Hjelm2020}

That is exactly where modern tools become useful. The right next question is not ``Do we believe the cipher?'' It is ``How often does a thing this good happen by accident?''

\end{JournalAbstract}
\section*{The river that tells you how to read it}

Genesis 2:10--14 is a strange little island in the Eden story. The narrative has just introduced the garden, its trees, the tree of life, and the tree of the knowledge of good and evil. Then it pauses for a river. The river waters Eden, leaves it, and ``from there'' divides into four heads. The first is Pishon, associated with Havilah, gold, bdellium, and the \textit{shoham} stone. The second is Gihon. The third is Hiddekel, usually identified with the Tigris in translation. The fourth is Euphrates. Then, almost as abruptly, the narrator goes back to the human being in the garden.\cite{MasoreticGenesis}

For readers without Hebrew, here is the machinery in Genesis 2:10 that matters. There is no need to memorize it.

\JournalTable{|p{0.18\linewidth}|p{0.24\linewidth}|p{0.48\linewidth}|}{%
\textbf{Hebrew} & \textbf{Transliteration} & \textbf{Plain English} \\ \hline
\hebrew{ומשם} & \textit{u-misham} & and from there \\ \hline
\hebrew{יפרד} & \textit{yippared} & it divides / branches apart \\ \hline
\hebrew{והיה} & \textit{ve-hayah} & and it becomes \\ \hline
\hebrew{לארבעה} & \textit{le-arba'ah} & into four \\ \hline
\hebrew{ראשים} & \textit{roshim} & heads / headwaters \\ \hline
}

The Masoretic text really does move from ``from there'' to ``it divides'' to ``four heads''; this is not an interpretive paraphrase built for the present theory.\cite{MasoreticGenesis}

That sequence gives us a useful metaphor but not yet evidence. A river that divides is not proof that its paragraph is a cryptogram. Still, it tells us where to look: \textbf{division, location, naming, and repeated boundaries are already what the paragraph is about.}

There is another small structural fact worth filing away. The first three rivers are introduced with the ordinary Hebrew noun \hebrew{שם} (\textit{shem}, ``name''), either directly or with the conjunction ``and.'' The fourth is not. Genesis 2:14 simply says, in effect, ``and the fourth river is Euphrates.'' The prose is perfectly grammatical. There is no missing word that needs repair. But the asymmetry is real.\cite{MasoreticGenesis}

\JournalTable{|p{0.14\linewidth}|p{0.20\linewidth}|p{0.28\linewidth}|p{0.28\linewidth}|}{%
\textbf{River} & \textbf{Hebrew} & \textbf{Transliteration} & \textbf{What it does} \\ \hline
Pishon & \hebrew{שם} & \textit{shem} & name-marker \\ \hline
Gihon & \hebrew{ושם} & \textit{u-shem} & ``and the name'' \\ \hline
Hiddekel & \hebrew{ושם} & \textit{u-shem} & ``and the name'' \\ \hline
Euphrates & \hebrew{והנהר} & \textit{ve-ha-nahar} & ``and the river,'' without the matching name-marker \\ \hline
}

Again: an omission is not a code. Prose varies. The point is simply that the paragraph makes the idea of a \textbf{name} visually repetitive and then breaks its own pattern on the fourth river.\cite{MasoreticGenesis}

The same three consonants used for ``name'' also occur inside the Pishon description with an entirely different ordinary reading: \hebrew{שם} (\textit{sham}, ``there''). Biblical Hebrew does this without strain. The same consonantal sequence can be a noun meaning ``name,'' an adverb meaning ``there,'' or---under a third vocalization---the past-tense verb \hebrew{שם} (\textit{sam}, ``he put/placed''). The placement verb is extremely common in Biblical Hebrew, and the locative ``there'' is likewise common.\cite{BDB}

That is our first real hinge.

To a reader of the pointed Masoretic text, \textit{shem}, \textit{sham}, and \textit{sam} are visibly and audibly different words because the vowels tell us which is intended. To an experiment conducted on the consonants alone, they are one three-letter packet with three ordinary Hebrew exits:

\textbf{name / there / placed.}

That would be mildly amusing anywhere. It is more interesting in Genesis 2, because the same chapter is already unusually saturated with both \textbf{placing} and \textbf{naming}. God ``places'' the human in the garden in Genesis 2:8; animals are brought to the human to see what he will call them in 2:19; and 2:20 explicitly describes the human giving names. The chapter's ordinary narrative vocabulary therefore gives all three readings of the packet a home before we have invented a secondary reading at all.\cite{MasoreticGenesis,BDB}

This is also a narrative that openly likes wordplay. Genesis 2:23 connects \textit{ish} (``man'') and \textit{ishah} (``woman'') by sound, and scholarship treats the phonetic relationship as obvious even while the later relation between Eve's name and ``living'' in Genesis 3:20 is considerably harder to explain historically. In other words, the text itself teaches its audience that names, sounds, and meanings may be allowed to rub against one another without becoming a modern dictionary exercise.\cite{Napora2022}

None of that proves our case. It does something more modest and more useful: it makes the proposed kind of reading \textbf{native to the literary neighborhood} rather than wholly alien to it.

A final local feature matters because it keeps us from making naïve claims about spelling. Genesis 2:11 and 2:13 use essentially the same expression, ``the one circling,'' of Pishon and Gihon, but one occurrence writes the long \textit{o} without a consonantal \textit{waw} while the other writes it with one. This is what scholars call \textbf{defective} versus \textbf{plene} spelling: roughly, a spelling with fewer versus more vowel-letters. Modern work emphasizes that the history of such spellings is complex; Joel Elitzur has specifically challenged the simple story in which an originally defective biblical text was mechanically filled out by later scribes, while Aaron Hornkohl's recent work likewise treats the Torah's orthographic profile as a serious historical phenomenon rather than random sloppiness.\cite{Elitzur2023,Hornkohl2024}

So the nearby variation does \textbf{not} prove microscopic authorial engineering. It tells us something safer: the consonantal form of this paragraph cannot be reasoned about as though modern standardized spelling rules were in force.

And one historical shortcut should be ruled out immediately. We should not say, ``Ancient Hebrew had no spaces, therefore later scribes guessed the word boundaries.'' A. R. Millard's classic study was directed precisely against routine appeals to \textit{scriptio continua}---continuous writing without word division---as an easy explanation for Hebrew resegmentation. Early Northwest Semitic scribal traditions had methods of marking word divisions. If the phenomenon proposed here is ancient, the interesting model is therefore \textbf{deliberate crossing of known boundaries}, not recovery of spaces that nobody had bothered to write.\cite{Millard1970}

That distinction is everything.

The proposal is not: \textit{the scribes got the sentence wrong.}

The proposal is: \textit{perhaps the sentence was built so that the right reading was not the only lexical path through its consonants.}

\section*{The hinge in the consonants}

Now we can put the cards face up.

The normal Masoretic reading at the end of Genesis 2:11 and through 2:12 is straightforward: in Havilah ``there'' is gold; the gold of that land is good; ``there'' are bdellium and the \textit{shoham} stone. Genesis 2:13 then begins ``and the name of the second river\textbackslash{}ldots{}.''\cite{MasoreticGenesis}

Here is the relevant surface text, deliberately broken into one Hebrew word per row.

\JournalTable{|p{0.14\linewidth}|p{0.20\linewidth}|p{0.28\linewidth}|p{0.28\linewidth}|}{%
\textbf{Verse} & \textbf{Hebrew} & \textbf{Transliteration} & \textbf{Ordinary value} \\ \hline
2:11 & \hebrew{שם} & \textit{sham} & there \\ \hline
2:11 & \hebrew{הזהב} & \textit{ha-zahav} & the gold \\ \hline
2:12 & \hebrew{וזהב} & \textit{u-zahav} & and the gold \\ \hline
2:12 & \hebrew{הארץ} & \textit{ha-aretz} & the land \\ \hline
2:12 & \hebrew{ההוא} & \textit{ha-hi} & that, feminine \\ \hline
2:12 & \hebrew{טוב} & \textit{tov} & good \\ \hline
2:12 & \hebrew{שם} & \textit{sham} & there \\ \hline
2:12 & \hebrew{הבדלח} & \textit{ha-bedolach} & the bdellium \\ \hline
2:12 & \hebrew{ואבן} & \textit{ve-even} & and a stone \\ \hline
2:12 & \hebrew{השהם} & \textit{ha-shoham} & the \textit{shoham} stone \\ \hline
2:13 & \hebrew{ושם} & \textit{u-shem} & and the name \\ \hline
}

The forms and their order are those of the Masoretic Genesis text.\cite{MasoreticGenesis}

For the experiment, strip the vowels and cantillation and ask only what happens if the spaces move. Also normalize the special final shape of \textit{mem}; final \hebrew{ם} and medial \hebrew{מ} are the same consonant /m/, and our comparison concerns consonants rather than positional typography.

In transliteration, the \textbf{surface segmentation} is:

\textit{sham | ha-zahav | u-zahav | ha-aretz | ha-hi | tov | sham | ha-bedolach | ve-even | ha-shoham | u-shem}

The \textbf{alternate segmentation} is:

\textit{sam | ha-zeh | buzah | bah | ertzeh | hu | tov | samah | b'dal | ḥwʾ | ben | ha-seh | musam}

After normalization, both segmentations consist of the same \textbf{39 consonants in the same order}. Nothing is inserted. Nothing is deleted. Nothing trades places. No \textit{dalet} becomes a \textit{resh}. No Paleo-Hebrew look-alike is invoked. The only operation is deciding where one lexical unit ends and the next begins.

That is what ``zero-edit'' means in this paper.

The final alternate word needs one brief explanation because it makes the experiment look longer than the user's stated target. The main field of play is exactly what matters most here: the final two words of 2:11 plus all of 2:12. At the end of 2:12, however, the alternate segmentation has carried an /m/ forward. The opening word of 2:13 immediately supplies the remaining consonants needed to resolve it. So 2:13 contributes a \textbf{one-word coda}, not a second major passage.

Here is the second cut.

\JournalTable{|p{0.14\linewidth}|p{0.20\linewidth}|p{0.28\linewidth}|p{0.28\linewidth}|}{%
\textbf{Hebrew} & \textbf{Suggested vocalization} & \textbf{Plain-English lexical value} & \textbf{Assessment} \\ \hline
\hebrew{שם} & \textit{sam} & he placed & very strong \\ \hline
\hebrew{הזה} & \textit{ha-zeh} & this / this one & very strong \\ \hline
\hebrew{בוזה} & \textit{buzah} & contempt / an object of contempt & strong \\ \hline
\hebrew{בה} & \textit{bah} & in her / with her & very strong \\ \hline
\hebrew{ארצה} & \textit{ertzeh} & I will be pleased / delight & very strong \\ \hline
\hebrew{הוא} & \textit{hu} & he / it & very strong \\ \hline
\hebrew{טוב} & \textit{tov} & good & very strong \\ \hline
\hebrew{שמה} & \textit{samah} & she placed & very strong as a form; context open \\ \hline
\hebrew{בדל} & \textit{b'dal} & a severed piece / fragment of & very strong \\ \hline
\hebrew{חוא} & \textit{ḥwʾ} & unresolved & weak \\ \hline
\hebrew{בן} & \textit{ben} & son / offspring & very strong \\ \hline
\hebrew{השה} & \textit{ha-seh} & the sheep-or-goat / flock animal & very strong \\ \hline
\hebrew{מושם} & \textit{musam} & placed / set & morphologically strong; register less secure \\ \hline
}

Several of these assessments can be checked against ordinary Biblical Hebrew lexicography rather than ingenuity. The noun represented by \textit{buzah} is attested with the meaning ``contempt'' or ``object of contempt'' in Nehemiah 4:4. The verb underlying \textit{ertzeh} means ``be pleased with'' or ``accept favorably''; Brown--Driver--Briggs explicitly records uses with the preposition \textit{b-} before both persons and things. The noun \textit{b'dal} is not merely a convenient root abstraction: it is an attested Biblical Hebrew noun meaning ``piece, severed piece,'' found in Amos 3:12 in the expression ``a piece of an ear.'' And \textit{seh} denotes one member of a small-livestock flock, a sheep or a goat; translating it automatically as ``lamb'' would be too specific.\cite{BDB}

Already we can see why the proposal deserves something better than either enthusiasm or ridicule. The alternate cut is not a row of thirteen invented noises. Almost the entire thing lands on recognizable Hebrew material.

But recognizable words are cheap.

A Scrabble rack full of English words is not an English sentence.

The question is what happens when neighboring pieces begin to constrain one another.

The strongest example is the pair represented by \textit{bah | ertzeh}. The first means ``in her'' or ``with her.'' The second can be the first-person form ``I will be pleased / I will delight.'' Since the Biblical Hebrew verb is actually attested with \textit{b-} complements, the two pieces do not merely sit beside one another as dictionary possibilities. They form a genuine governing relationship: approximately \textbf{``with her I will be pleased''} or \textbf{``in her I will delight.''}\cite{BDB}

This matters enormously. In pattern hunting, adjacent words that independently exist are one thing; adjacent words that discover the grammar they need in each other are better evidence.

The next pair is almost comically simple: \hebrew{הוא} (\textit{hu}, ``he/it'') followed by \hebrew{טוב} (\textit{tov}, ``good'') gives \textbf{``he/it is good.''}

There is, however, a delightful orthographic wrinkle here. In the ordinary surface sentence, these consonants are part of the expression translated ``that land.'' The noun for ``land'' is feminine, and the Masoretic reading accordingly treats the pronoun as feminine---\textit{ha-hi}, ``that [feminine].'' Yet the Tiberian Torah has the well-known peculiarity of often writing the third-person feminine singular pronoun with the consonants normally associated elsewhere with the masculine form. Hornkohl calls the Pentateuchal written tradition a notable outlier in exactly this respect. Genesis 2:12 therefore contains, in the surface text itself, the consonants that an alternate reading can hear as masculine \textit{hu}, ``he,'' while the ordinary reading correctly understands the form as feminine.\cite{Hornkohl2024,MasoreticGenesis}

This is precisely the kind of ambiguity worth paying for because the text already pays for it in the ordinary reading. We do not have to invent an exotic spelling rule for the secondary one.

The next major piece is \hebrew{בדל} (\textit{b'dal}, ``a severed piece''). This was the discovery that most improved the earlier version of the argument. The word need not be treated as ``something vaguely derived from the root meaning separation.'' The noun itself exists. Amos 3:12 uses it in construct form---\textbf{construct} simply means the Hebrew relationship corresponding roughly to English ``X of Y''---for a severed piece of an ear.\cite{BDB}

That gives the alternate segmentation a small but striking semantic rhyme with Genesis 2:10. The surface river ``divides'' using a different Hebrew root; the alternate segmentation later yields a noun meaning ``severed piece.'' These are \textbf{not cognates}, and no hidden etymology should be claimed. They simply occupy the same semantic neighborhood: division on the surface, a severed piece underneath.\cite{MasoreticGenesis,BDB}

Then comes \hebrew{בן} (\textit{ben}, ``son/offspring''), among the least controversial words in the whole experiment, followed by \hebrew{השה} (\textit{ha-seh}, ``the sheep-or-goat,'' more generically a flock animal). The latter is worth translating carefully because English Bible tradition can make \textit{seh} sound automatically like ``lamb.'' Lexically, the animal can be a sheep or goat; Exodus 12 itself makes that range explicit in context.\cite{BDB}

The last piece, \hebrew{מושם} (\textit{musam}, ``placed/set''), is morphologically transparent as a passive ``placed'' formation from the extraordinarily common Hebrew verb meaning ``put/place/set.'' What should make us cautious is not whether Hebrew can form such a word, but how much weight we should place on this exact form as representative of the biblical register. The root itself is deeply biblical and appears in Genesis 2:8 in the act of placing the human in Eden; the exact secondary form should therefore be described as morphologically good without pretending its corpus credentials are as clean as those of \textit{ben} or \textit{b'dal}.\cite{BDB}

And now we hit the wall.

\hebrew{חוא} (\textit{ḥwʾ}, here left untranslated) is not ordinary Biblical Hebrew for Eve. The biblical name is \hebrew{חוה} (\textit{Chavvah}, ``Eve''), with a different final consonant. Genesis 3's association between Eve and ``living'' is itself philologically complicated; modern scholarship has surveyed a long history of attempts to explain it. That context makes life/Eve comparisons interesting, but it does not license changing an \textit{aleph} into a \textit{he}.\cite{Napora2022}

Earlier exploratory work can find comparative Semitic material that makes \textit{ḥwʾ} less than meaningless, especially in later Aramaic/Syriac territory. But this revision deliberately refuses to let comparative languages perform a rescue operation that Biblical Hebrew has not earned.

That single decision improves the paper.

The secondary stream should \textbf{not} currently be printed as a complete hidden English sentence.

It is better represented as a sequence of locally strong clauses with one unresolved packet:

\JournalQuote{\textit{sam ha-zeh buzah} --- perhaps ``this one placed contempt,'' though compressed and not load-bearing.}

\JournalQuote{\textit{bah ertzeh} --- ``with/in her I will be pleased; I will delight.''}

\JournalQuote{\textit{hu tov} --- ``he/it is good.''}

\JournalQuote{\textit{samah b'dal \textbackslash{}ldots{} ben} --- ``she placed a severed piece of \textbackslash{}ldots{} a son/offspring,'' with the crucial complement unresolved.}

\JournalQuote{\textit{ha-seh musam} --- ``the flock animal is placed/set.''}

The first of those is plausible enough to preserve but not smooth enough to make foundational. The middle pair is substantially better because the verb \textit{ratzah} genuinely takes the prepositional construction at issue. The ``severed piece'' noun is exact Biblical Hebrew. The final animal-placement pair is morphologically natural but should carry less evidential weight because of the precise register question around \textit{musam}.\cite{BDB}

The result is a funny kind of success. We started wanting a secret sentence and found something intellectually healthier: a \textbf{strong lexical architecture with an exposed bad joint}.

A theory with no bad joints is usually a theory whose author has hidden them.

\section*{What the alternate words actually buy us}

The most interesting thing in the passage is not any one translated clause. It is the way several mechanisms stack.

Start with the little three-letter machine \hebrew{שם} (\textit{shem/sham/sam}, ``name/there/he placed''). On the surface of Genesis 2:11--14, the packet oscillates between \textbf{name} and \textbf{there}. Under the alternate vocalization, it can move again into \textbf{placed}. The larger Eden narrative independently uses the placement verb for God putting the human in the garden and foregrounds the act of naming animals shortly afterward.\cite{MasoreticGenesis,BDB}

Then the alternate segmentation produces \hebrew{שמה} (\textit{samah}, ``she placed''), an exact ordinary third-person feminine form of the same placement verb; BDB explicitly lists that form and its biblical occurrence.\cite{BDB}

And at the far end we get the passive-looking \textit{musam}, ``placed.''

So one possible structural frame is:

\textbf{he placed → she placed → placed.}

At the same time, the first of those three consonantal families continues to hover over the ordinary meanings:

\textbf{name → there → place.}

That is a much stronger observation than ``look, some Hebrew words appeared.'' It is a repeated morphological idea produced at separated points in one continuous zero-edit resegmentation.

It also resonates with what Genesis 2 is doing narratively. The chapter is full of things being situated: a garden is planted, a human is placed, trees are made to grow, a river issues and divides, regions are described, animals are brought, and names are assigned. This is a story obsessively interested in the act of putting things where they belong and saying what they are.\cite{MasoreticGenesis}

The second striking feature is the \textit{bah | ertzeh} pair. It is tempting to rush straight to theology---``in her I delight''---but philologically the more important point comes first. The proposed segmentation creates a prepositional complement and then immediately creates a verb known to govern that kind of complement. That is the sort of coincidence a computational test ought to reward more heavily than two unrelated dictionary hits.\cite{BDB}

The third feature is the surface/secondary pronoun switch in Genesis 2:12. The ordinary reading requires the consonants to function as feminine ``that,'' agreeing with the feminine noun ``land''; the alternate boundary frees the inner packet as masculine \textit{hu}, ``he.'' Because the Pentateuch's written tradition really does exhibit this masculine-looking spelling for the feminine pronoun, we are not manufacturing the ambiguity. It belongs to the received text.\cite{Hornkohl2024,MasoreticGenesis}

Then comes the rare noun \textit{b'dal}, ``severed piece.'' Rarity can be abused in two opposite ways. A rare word is not automatically exciting: large texts contain rare things. But when a resegmentation lands on an \textbf{actually attested rare lexeme} with an unusually apt local meaning, that should count for more than landing on a common preposition. In Amos 3:12 the noun denotes a piece torn or saved from an ear; in Genesis 2 the surface paragraph has just defined its river by division. Again, the roots differ, so this is semantic echo, not etymological identity.\cite{BDB,MasoreticGenesis}

Finally there is the cluster ``son/offspring'' and ``flock animal.'' Those words tempt a reader who knows what happens later in Genesis to sprint toward sacrifice, sons, substitutions, Abel, Isaac, Passover, or Christ. That is precisely where discipline matters.

At this stage, those later connections belong in a search queue, not in the conclusion.

A pattern becomes less impressive, not more, when its interpreter is permitted to draw an arrow from every evocative noun to every later biblical story.

The right question is narrower: \textbf{does the secondary segmentation produce an unusually dense sequence of valid forms and locally coherent grammatical relationships?}

At present, I think the answer is yes.

Does that mean an ancient writer intended it?

Not yet.

There is a useful asymmetry in the evidence:

\JournalTable{|p{0.18\linewidth}|p{0.24\linewidth}|p{0.48\linewidth}|}{%
\textbf{Feature} & \textbf{What we actually have} & \textbf{Weight} \\ \hline
Same consonants & exact zero-edit identity after normalizing positional final forms & high \\ \hline
\textit{bah + ertzeh} & a real complement + governing verb relationship & high \\ \hline
\textit{hu + tov} & elementary Hebrew clause & high \\ \hline
\textit{b'dal} & exact biblical noun, ``severed piece'' & high \\ \hline
\textit{ben} & ordinary ``son/offspring'' & high \\ \hline
\textit{ha-seh} & ordinary flock-animal noun & high \\ \hline
placement pattern & several forms cluster around ``put/place'' & moderate to high \\ \hline
``there/name'' oscillation & overtly present in the surface paragraph & moderate to high \\ \hline
river/name asymmetry & first three named with marker, fourth without it & supporting \\ \hline
\textit{buzah} opening & lexical, but syntax remains compressed & moderate \\ \hline
\textit{musam} & transparent morphology, less secure exact biblical register & moderate \\ \hline
\textit{ḥwʾ} & no secure ordinary Biblical Hebrew solution & low \\ \hline
}

The lexical values in this summary follow the Masoretic text and standard lexical evidence discussed above.\cite{MasoreticGenesis,BDB}

Notice what is \textit{not} in the table: ``Phoenician secretly says life,'' ``Paleo-Hebrew makes two letters interchangeable,'' ``the hidden text says Eve,'' ``there are four secret messages because there are four rivers.''

Those are exactly the sorts of moves that make an intriguing observation impossible to falsify.

Earlier versions of this argument tried to give names---Stream A, Stream B, Stream C---to several overlapping possibilities. I now think that was too generous. A theory gets stronger when it is deprived of escape routes.

There is one main alternate segmentation.

There are several local ways of vocalizing ambiguous pieces.

There is one serious lexical bottleneck.

And there is one empirical question: \textbf{is the whole configuration more structured than comparable accidents elsewhere?}

That is enough.

One additional fact makes the \textit{hu} portion of the pattern especially worth testing but also makes the manuscript problem unavoidable. The Samaritan Pentateuch's Genesis 2:12 does not have precisely the Masoretic consonantal string. It uses the ordinary feminine spelling corresponding to \textit{hi} rather than the Torah's Masoretic \textit{hwʾ} spelling, and it adds ``very'' after ``good.'' Either change disrupts the proposed carry-chain; together they break it decisively.\cite{SamaritanGenesis}

The Samaritan Pentateuch should not be dismissed as a late scribal curiosity just because it spoils our toy. It is an important independent Pentateuchal textual tradition, and current scholarship resists treating it as merely an aberrant derivative of the Masoretic form. A 2024 scholarly dataset project likewise describes the Samaritan Pentateuch as an important early witness and now makes its text available with word-level linguistic annotation.\cite{Hjelm2020,Naaijer2024}

That means the right historical claim is narrower than the exciting one:

\JournalQuote{\textbf{The resegmentation is demonstrably present in the consonants of the Masoretic textual tradition represented by our medieval witnesses.}}

We do \textbf{not} yet get to replace that sentence with:

\JournalQuote{\textit{The author of Genesis designed it.}}

That second sentence requires a history of the reading, not merely a clever parsing.

The Dead Sea Scrolls cannot currently settle this particular point for us. The surviving material conventionally catalogued as 4Q2 preserves Genesis 2:14--19 after earlier material in Genesis 1, but not 2:11--13. So the crucial consonantal window is absent from that witness.\cite{Tov2012}

This is frustrating, but it is exactly the kind of frustration textual criticism exists to enforce.

A missing ancient witness is not an invitation to become more certain.

It is an empty chair at the table.

\section*{Where the argument can fail}

Any paper of this kind needs an adversary. Fortunately we have one: Hebrew.

Hebrew is an excellent machine for manufacturing accidental patterns. It has short roots. It attaches conjunctions, articles, prepositions, and pronouns directly to words. An unvocalized string can often support more than one reading. Some common consonantal packets correspond to multiple unrelated lexemes. If one is allowed to slide boundaries freely and then hunt through Hebrew, Aramaic, Phoenician, Ugaritic, Akkadian, Arabic, Syriac, and modern Hebrew until something interesting appears, the probability of ``discovery'' approaches one.

So the first rule is brutally simple:

\textbf{A multilingual rescue is a penalty, not a bonus.}

Comparative Semitic languages are indispensable for historical linguistics. They can illuminate cognates, older meanings, and possible wordplay. But a form that is ordinary Biblical Hebrew is better evidence for a Biblical Hebrew secondary reading than a form that becomes meaningful only after crossing language and century boundaries.

That is why \textit{ḥwʾ} remains unresolved here.

The second rule is:

\textbf{Lexicality is not syntax.}

``Lexicality'' simply means that a chunk is a genuine word or word-form. Syntax is the machinery by which words make clauses. Thirteen genuine words do not automatically make one genuine sentence. The strongest parts of this experiment are the places where syntax begins to emerge without special pleading---especially \textit{bah ertzeh}, ``with/in her I will be pleased,'' and the trivial but clean \textit{hu tov}, ``he/it is good.''\cite{BDB}

The third rule is:

\textbf{Do not change a consonant.}

There are legitimate cases in textual criticism where graphic similarity between letters helps explain variants, and recent work on differences between Masoretic and Samaritan traditions continues to investigate exactly such paleographic processes. But that is a different operation from the one being tested here. Our present pattern has an unusual virtue: it does not need substitutions. The moment we start turning one consonant into another because a prettier hidden sentence appears, we have changed experiments.\cite{Tov2012}

The fourth rule is:

\textbf{Do not pretend the original manuscript was a strip of letters with no word boundaries.}

As noted above, ancient Hebrew and related Northwest Semitic writing traditions had word-dividing practices. Intentional boundary play remains perfectly conceivable---English writers can pun across spaces even though English has spaces---but ``the ancients did not know where the words were'' is not a sound historical foundation.\cite{Millard1970}

The fifth rule is the cruelest:

\textbf{The Samaritan variant counts against us.}

It is tempting to describe a variant that destroys a theory as corruption and a variant that preserves it as antiquity. That is not textual criticism; that is gardening.

The Samaritan reading shows that the precise consonantal machine under discussion was not invariant across the Pentateuch's transmission. At minimum, one of three broad possibilities follows: the Masoretic-type sequence preserves an old design that the Samaritan branch disturbed; the design arose or became exact only within the textual history ancestral to the Masoretic branch; or the whole effect is accidental and the variant merely exposes how contingent it is. Present evidence does not select among those possibilities.\cite{SamaritanGenesis,Naaijer2024,Hjelm2020}

The sixth rule is:

\textbf{The normal text always gets first rights.}

The ordinary Genesis 2:11--14 reading is not broken Hebrew waiting for this paper to rescue it. It is coherent narrative prose. It names rivers and describes places. The identities of some materials and geographical referents may be debated, but the sentence structure at issue is not in need of our resegmentation.\cite{MasoreticGenesis}

That is an important distinction between \textbf{replacement readings} and \textbf{multiplexed readings}.

A replacement reading says, ``Everyone has divided this incorrectly.''

A multiplexed reading says, ``The obvious reading may be intentionally complete while another pattern is also latent in the same material.''

Only the second hypothesis is being defended as worth testing.

And there is one final danger: \textbf{semantic magnetism}.

Once a reader knows that Genesis contains Eve, sons, flock animals, sacrifice, death, life, brothers, blood, substitution, and later covenantal narratives, ordinary words begin to behave like iron filings around a magnet. ``Son'' suddenly points to Isaac. ``Flock animal'' points to Passover. ``Good'' points to creation. ``Severed piece'' points to the woman's creation. ``Placed'' points to Eden. Soon every word points everywhere.

That feeling of convergence can be psychologically powerful while carrying almost no statistical information.

The cure is not less imagination.

The cure is controls.

\section*{How to test this instead of believing it}

This is the part Newton did not have.

Newton's fascination with Scripture and prophecy was not a minor eccentric footnote to his intellectual life. The Newton Project has published enormous quantities of his theological, prophetic, chronological, and church-historical writing; its account of his work describes more than two million surviving words on early church history alone, alongside extensive efforts on biblical prophecy, Jewish worship, chronology, and textual variants. The Project's catalogue includes a roughly 76,500-word draft work on Daniel and Revelation, and its larger digital edition has made most of the surviving religious papers available for study.\cite{NewtonProject}

Newton had the obsession.

We have the corpus.

That changes the proper standard of argument.

A reader in Newton's position could find a pattern, compare passages by hand, consult available languages, and argue that the resulting architecture looked intentional. We can still do all those things. But we can now ask the one question that hand scholarship is terrible at answering:

\textbf{Compared with what?}

A \textbf{null model} is simply a model of what accident looks like.

Suppose I tell you I flipped ten coins and obtained seven heads. Is that remarkable? You cannot answer by staring harder at the sequence. You need to know how often seven or more heads occurs under fair flipping.

The same principle applies here.

Finding \textit{b'dal}, ``severed piece,'' feels remarkable. Finding \textit{bah ertzeh}, ``with her I will be pleased,'' feels more remarkable. Finding both in one zero-edit segmentation feels more remarkable still.

But until we know how often comparable Hebrew strings generate comparable alternate segmentations, ``remarkable'' is autobiography.

Fortunately, a large part of the necessary infrastructure already exists. The Eep Talstra Centre for Bible and Computer maintains a Hebrew Bible database with linguistic annotation at word, phrase, clause, and textual levels; its BHSA data are available for computational research through Text-Fabric, while SHEBANQ provides queryable access to the annotated corpus. This is not speculative future infrastructure. Biblical Hebrew has been computationally encoded for decades.\cite{ETCBC}

The broader field has now reached the point where a 2025 chapter in \textit{The Oxford Handbook of Textual Criticism of the Bible} describes a ``technological turn'' involving manuscript databases, material science, digital imaging, computational methods, electronic editions, and digital research environments. The argument of this paper therefore does not require biblical studies to become computational. That process is already under way.\cite{Longacre2025}

For this particular problem, the experiment should be simple enough to explain to a bright teenager.

First, freeze the rules \textbf{before} doing the large search.

A valid alternate segmentation gets the same normalized consonants in the same order. No substitutions. Each candidate word receives a score based on how securely it belongs to Biblical Hebrew. Ordinary corpus-attested forms score highly. Morphologically regular but unattested or later forms score less. A jump into another Semitic language incurs a large penalty. Grammatical relationships between adjacent chunks add evidence. Semantic resemblance to the local passage can be scored separately so that ``these are real words'' never gets confused with ``these words are thematically interesting.''

Second, run the same procedure on lots of biblical windows of comparable length.

Not just famous passages. Not just Genesis. Random 39-consonant stretches from narrative, law, poetry, and prophecy.

Third, create harder controls.

Shuffle consonants. Preserve individual letter frequencies. Then preserve two-letter and three-letter frequencies so the fake strings still resemble Hebrew. Generate controls that maintain common prefixes. Generate controls whose morphology looks Hebrew but whose local semantics are nonsense.

The more Hebrew-like the control becomes, the more interesting Genesis must remain to survive.

Fourth, measure the thing that human analysts usually hide: \textbf{degrees of freedom}.

How many starting points did we try before finding this one?

How many alternate vocalizations were considered?

How many dictionaries?

How many languages?

How many times was a boundary adjusted after seeing an attractive downstream word?

How many semantic interpretations were entertained and discarded?

That ledger is part of the evidence. An astonishing pattern found under one frozen rule is different from the same pattern found after a thousand invisible nudges.

Fifth, split discovery from confirmation.

The present Genesis passage has already been looked at too much to serve as pristine test data. Fine. Use it to build the scoring system, then lock the scoring system and send it blind into held-out biblical text. If it rediscovers passages already independently recognized for overt wordplay, that is encouraging. If it begins flagging arbitrary prose with scores as high as Genesis 2:11--12, our miracle has evaporated.

The machine-learning analogy is useful here precisely because machine learning has learned this lesson the hard way: a model that performs brilliantly on the material used to tune it may simply have memorized the analyst.

Sixth, make the results reproducible.

A genuine computational argument should eventually allow another group to download the corpus, run the same segmentation code, use the same lexical inventory, see the same candidate rankings, change the penalties, and discover exactly when the result disappears. ETCBC's data infrastructure was built in part around precisely this ideal of shared, queryable, reusable scholarship.\cite{ETCBC}

AI belongs here, but in a less glamorous role than people usually assign it.

A language model should \textbf{not} be asked, ``Was Genesis secretly engineered this way?'' and then cited because it answered confidently.

That is astrology with GPUs.

Modern models are useful as laboratory equipment: generating candidate segmentations, proposing vocalizations, ranking semantic coherence, searching cognate databases, finding remote parallels, producing synthetic controls, and explaining candidate parses to non-specialists. But every one of those outputs can be audited against a corpus, lexicon, grammar, manuscript, or explicit scoring function.

The distinction is simple:

\textbf{AI may accelerate hypothesis generation. It may not manufacture evidence by having an opinion.}

This is where the modern project can become genuinely better than both old confessional exegesis and old academic skepticism.

The believer's temptation is to say, ``The pattern is beautiful; therefore God put it there.''

The skeptic's temptation is to say, ``People find patterns; therefore there is nothing to investigate.''

Both responses skip the experiment.

The twenty-first-century response should be:

\textbf{freeze it, count it, compare it, break it.}

And if it refuses to break, then become more interested.

\section*{Exile from Ivory, toward Babel}

What would a proper modern biblical sect look like?

``Proper'' is doing a great deal of work there.

I do not mean a congregation organized around the conviction that this paper has decoded Eden. That would be an unusually efficient way to ruin both religion and philology.

I mean something closer to what Newton wanted and could not quite build: a community that treats Scripture as important enough to submit every claim about it to the strongest tools human beings possess.

Newton's religious project was restorationist. He thought centuries of institutions could accumulate error and that ancient sources had to be reopened rather than merely inherited through official formulations. Whatever one thinks of his actual theological conclusions, the intellectual instinct is recognizable: go back to the evidence, distinguish inherited authority from demonstrated fact, and refuse to regard a venerable consensus as self-authenticating.\cite{NewtonProject}

The modern version should be much less Newtonian about Newton himself.

No solitary genius gets to be pope.

No model gets to be oracle.

No scholar gets tenure over reality.

No doctrine gets immunity from manuscript evidence.

No skeptical convention gets immunity either.

The center would be the Bible in the historically richest sense of ``Bible'' we can now manage: not one English translation floating above history, and not even one Hebrew codex treated as though transmission never happened, but the whole astonishing object---Masoretic manuscripts, Samaritan Torah, Dead Sea Scrolls, Septuagint, ancient versions, scribal practices, archaeology, comparative Semitic languages, reception histories, linguistic databases, material analysis, and the accumulated record of people arguing about all of it. Contemporary textual criticism increasingly treats this multiplicity as the field itself rather than as noise to be cleared away on the path to one frictionless ``original text.''\cite{Longacre2025,Hjelm2020}

A modern sect built around that Bible would have to love uncertainty almost as much as revelation.

Its creed would need version numbers.

Not because truth changes, but because \textbf{our map of truth does}.

``Current confidence: 0.72'' is sometimes a more reverent sentence than ``Thus saith the Lord.''

Its scholars would publish negative results. A beautiful proposed acrostic that fails corpus controls would be celebrated as progress. An archaeological date that destroys a favorite interpretation would be good news. A Samaritan reading that breaks our clever Genesis segmentation would go in the main text, not in a footnote designed to make it disappear.

Its liturgy, if it had one, might include the phrase:

\JournalQuote{\textit{What evidence would change our minds?}}

That sounds sterile until one notices how close it is to repentance.

The community would also need multiple kinds of people. The Hebrew expert who cannot code is useful. The statistician who cannot parse Biblical Hebrew is useful. The manuscript conservator is useful. The archaeologist is useful. The rabbi, priest, pastor, atheist historian, literary critic, and suspicious undergraduate are all useful because each notices a different way a claim can go wrong.

That is not relativism.

It is error correction.

The old academic Ivory Tower has failed, and in every sense that matters for this project it has now ceased to exist. The important knowledge is no longer containable inside a few departments, a few journals, a few expensive monographs, and a credentialed priesthood that receives permission to speak in twenty-minute conference slots. The present technological turn in biblical textual criticism itself points in the opposite direction: databases, digital editions, computational tools, interoperable corpora, high-resolution images, and interdisciplinary workflows are pushing scholarship toward distributed technical collaboration.\cite{Longacre2025,ETCBC}

Good riddance to the monopoly. Not good riddance to expertise.

That distinction matters.

The failure of a gate does not imply the uselessness of the people who learned what was behind it. Biblical Hebrew still takes years. Paleography still takes years. Statistics still takes years. A person with a browser and a language model has not become Emanuel Tov by Tuesday afternoon.

The post-academic answer is therefore not anti-expert populism. It is \textbf{expertise without feudalism}.

Everything inspectable.

Everything challengeable.

Everything citeable.

Everything forkable where copyright and preservation allow it.

Reputation matters, but it never outranks the object under discussion.

In that sense, perhaps ``tower'' is still the right metaphor, just not Ivory.

We are in exile from Ivory.

We are moving toward Babel.

That sounds like a warning because Genesis 11 taught the Western imagination to hear Babel as failure: one human project, one ascending tower, one attempt to make a name, then fractured speech. Yet a modern Babel can invert the old architecture. It need not be a single hierarchy trying to reach heaven. It can be a place where many languages, disciplines, textual traditions, machines, and human communities meet without first pretending they are one language.

Hebrew beside Aramaic.

Masoretic text beside Samaritan.

A medieval codex beside a Qumran fragment.

Philology beside statistics.

Radiocarbon beside paleography.

Python beside prayer.

A rabbinic reading beside a Christian one beside an atheist one, with nobody permitted to win by changing the evidence.

That Babel would be noisy.

Good.

Noise is what a truth-seeking system sounds like before error correction.

The strange resegmentation in Genesis 2 offers a miniature example of the ethic. We found something genuinely interesting. We made it better by learning more Hebrew. Then we made it better again by refusing to translate one word. Then an independent textual tradition arrived and damaged the historical claim. Instead of being an embarrassment, that damage told us what the claim actually is.

That is the whole game.

At present we can say something precise.

In the Masoretic consonantal tradition of Genesis 2:11b--12, with the opening word of 2:13 serving as a coda, the letters support a second zero-edit segmentation containing an unexpectedly high concentration of recognizable Hebrew forms. Several adjacent units generate real grammatical relationships. Some resulting meanings---place, goodness, severance, offspring---echo the conceptual world around them. One packet remains seriously unresolved. The Samaritan Pentateuch does not preserve the consonantal conditions required for the same pattern. No surviving Dead Sea Scroll witness presently gives us those target verses to arbitrate the difference.\cite{MasoreticGenesis,BDB,SamaritanGenesis,Tov2012}

That is not a decoded revelation.

It is better.

It is a sharp question.

Genesis says that a river goes out from Eden, and from there it divides.

The ordinary text gives us four rivers.

The consonants may give us another division.

Now that we have the tools Newton lacked, belief is the least interesting thing we can do with it.

We can measure.

\begin{thebibliography}{99}

\bibitem{BDB}
Francis Brown, S. R. Driver, and Charles A. Briggs.
\newblock \textit{A Hebrew and English Lexicon of the Old Testament}.
\newblock Oxford: Clarendon Press, 1907.

\bibitem{Elitzur2023}
Joel Elitzur.
\newblock ``Plene Spelling and Defective Spelling in the Hebrew Bible: The Question of Dating.''
\newblock \textit{Journal of the American Oriental Society} 143, no. 4 (2023): 745--765.

\bibitem{ETCBC}
Eep Talstra Centre for Bible and Computer.
\newblock \textit{Biblia Hebraica Stuttgartensia Amstelodamensis / BHSA} and SHEBANQ resources.
\newblock Amsterdam: Vrije Universiteit Amsterdam.

\bibitem{Hjelm2020}
Ingrid Hjelm.
\newblock ``The Samaritan and Jewish Versions of the Pentateuch: A Survey.''
\newblock \textit{Religions} 11, no. 2 (2020): 85.

\bibitem{Hornkohl2024}
Aaron D. Hornkohl.
\newblock \textit{Diachronic Diversity in Classical Biblical Hebrew}.
\newblock Cambridge: Open Book Publishers, 2024.

\bibitem{Longacre2025}
Drew Longacre.
\newblock ``The Practice of Hebrew Bible Textual Criticism in the Twenty-First Century.''
\newblock In \textit{The Oxford Handbook of Textual Criticism of the Bible}, edited by Sidnie White Crawford and Tommy Wasserman.
\newblock Oxford: Oxford University Press, published online 2025.

\bibitem{MasoreticGenesis}
\newblock Masoretic Genesis text.
\newblock Genesis 2:10--14, including the Pishon description, the river-name sequence, and the relevant locative uses.

\bibitem{Millard1970}
A. R. Millard.
\newblock ```Scriptio Continua' in Early Hebrew: Ancient Practice or Modern Surmise?''
\newblock \textit{Journal of Semitic Studies} 15, no. 1 (1970): 2--15.

\bibitem{Naaijer2024}
Martijn Naaijer, Christian Canu Højgaard, Stefan Schorch, and Martin Ehrensvärd.
\newblock ``Text-Fabric Dataset of the Samaritan Pentateuch.''
\newblock \textit{Research Data Journal for the Humanities and Social Sciences} 9, no. 1 (2024).

\bibitem{Napora2022}
Krzysztof Napora.
\newblock ```Snake', `Life' or `Mother of All Living'? The Meaning of the Name Eve and Its Role in Gen 3.''
\newblock \textit{Verbum Vitae} 40, no. 1 (2022): 37--50.

\bibitem{NewtonProject}
Newton Project, University of Oxford.
\newblock Digital edition of Isaac Newton's theological, prophetic, chronological, and scientific manuscripts.

\bibitem{SamaritanGenesis}
\newblock Samaritan Pentateuch, Genesis 2:12.
\newblock Witness to a materially different consonantal condition from the Masoretic sequence tested here.

\bibitem{Tov2012}
Emanuel Tov.
\newblock \textit{Textual Criticism of the Hebrew Bible}.
\newblock 3rd ed. Minneapolis: Fortress Press, 2012.

\end{thebibliography}

\end{document}
